% Study sheet: Volume III, Chapter 11, Durable Execution.
% Numbers: \val{...} from durable_execution_studysheet_numbers.tex, written by durable_execution_studysheet_numbers.py.
\documentclass{studysheet}
\input{durable_execution_studysheet_numbers}
\sheetvolume{3}
\sheetchapter{11}{Durable Execution}
\sheetsources{Numbers from the chapter's mlsysim scenarios (hypothetical); questions marked assumed use new inputs.}

\begin{document}
% ===================================================================== FRONT: LEARN
\frontside
\begin{bigidea}
Long trajectories outlive their workers. The runtime keeps an \hl{append-only log} on storage that
survives a declared failure, writes each \hl{intent before dispatch}, and rebuilds state by folding
it, never by asking the model.
\end{bigidea}

\begin{leftcol}
\sheetsec{Key ideas}
\begin{ideas}
\idea{Durable against a named failure.}{An acknowledgment survives only the faults its storage
  survives: declare the failure domain first.}
\idea{The log is the record; state is a fold.}{A reducer replays typed events into the control record,
  calling no model or tool; the context is a view.}
\idea{Intent before dispatch.}{An admitted operation's grant and reservation are durable before the
  request leaves, or a crash leaves a phantom effect.}
\idea{Resume is authorization, not restart.}{Prove the prefix; take a new epoch the log, gate and resources check;
  restore and reconcile before admitting new work.}
\idea{In doubt is settled by evidence.}{Reconcile by operation identity; repeat only under an idempotency contract.}
\idea{Replay, not rerun; keep what the promise needs.}{A rerun samples a new trajectory. Retention
  follows the recovery promise; a digest keeps no bytes.}
\end{ideas}

\vspace{8pt}
\sheetsec{Three uses of a recording}
\begin{sheettable}{@{}l L@{}}
\thead{Mode} & \thead{What runs; what it supports}\\ \midrule
Time-travel & reducer folds events: the state held\\
Harness regression & harness on recorded inputs, effects off: how it responds\\
Counterfactual & changed condition, new observations: this experiment only\\
\end{sheettable}
\vspace{3pt}
{\fontsize{8.6}{11}\selectfont\color{muted}Divergence: from a verified checkpoint, the first mismatched step
boundary localizes; it does not prove cause.\par}

\vspace{8pt}
\fcard{Resume budget}%
  {\ResumeTime = \OwnershipAcquisitionTime + \CheckpointSuffixReadTime + \SuffixFoldTime
   + \WorkspaceRestoreTime + \OperationReconciliationTime + \PrefillTime}%
  {Ownership, checkpoint and suffix read, fold, workspace restore, reconciliation, prefill; overlapping stages count
   by critical path.}%
  {With no cached prefix, prefill adds a floor of about \val{coldprefill}~s for
   \val{coldtokens} tokens; folding calls no model.}{}
\end{leftcol}%
\begin{rightcol}
\sheetsec{Where the crash falls}
{\centering
\begin{tikzpicture}[sheet]
  \foreach \x/\lab/\sty in {21/proposal/nbox, 67/admission/nbox, 113/intent ack/accbox, 159/dispatch/nbox, 205/evidence/nbox}{
    \node[\sty, minimum width=38pt, minimum height=15pt] at (\x,60) {\lab};}
  \foreach \a/\b in {40/48, 86/94, 132/140, 178/186}{\draw[arr] (\a,60) -- (\b,60);}
  \draw[dim, dashed] (136,76) -- (136,22);
  \node[note] at (136,82) {durability barrier: the request may leave};
  \draw[dim] (113,51) -- (113,38); \draw[dim] (205,51) -- (205,38);
  \node[nbox, fill=rulegray!40, minimum width=111pt, minimum height=14pt] at (55.5,30) {not dispatched};
  \node[hotbox, minimum width=88pt, minimum height=14pt] at (159,30) {in doubt};
  \node[accbox, minimum width=43pt, minimum height=14pt] at (228.5,30) {established};
  \node[note, anchor=north] at (55.5,20) {intent may still be durable};
  \node[note, anchor=north] at (159,20) {reconcile; hold dependents};
  \node[note, anchor=north] at (228.5,20) {per endpoint};
\end{tikzpicture}\par}
\vspace{4pt}
{\fontsize{8.6}{11}\selectfont A found intent cannot show whether the request left. A missing one
proves non-dispatch only in a \hl{complete authoritative prefix}, never a lagging replica.\par}

\vspace{8pt}
\sheetsec{What each acknowledgment survives}
\begin{sheettable}{@{}l L@{}}
\thead{Failure} & \thead{Acknowledgment that survives}\\ \midrule
Process crash & returned \texttt{write()} (page cache)\\
Power or host loss & \texttt{fdatasync()}, drive cache flushed\\
Volume loss & quorum of replicas off that volume\\
\end{sheettable}

\vspace{8pt}
\fcard{Checkpoint stride}%
  {\begin{gathered}\CheckpointRecoveryOverheadTime = \StrideCheckpointCount\,\CheckpointWriteTime
   + \TraceCrashCount\,\CheckpointRestoreTime + \TurnReconstructionTime \cdot \text{lost turns}\\
   \CheckpointStride^* = \sqrt{2\,\CheckpointWriteTime/(\TurnFailureRate\,\TurnReconstructionTime)}\end{gathered}}%
  {Every $\CheckpointStride$ turns: $\StrideCheckpointCount$ pauses of $\CheckpointWriteTime$. Each of
   $\TraceCrashCount$ crashes restores ($\CheckpointRestoreTime$), then rebuilds each turn past the last checkpoint ($\TurnReconstructionTime$ each). $\TurnFailureRate$: crashes $\div$ turns.}%
  {Unlogged turns cost re-execution $\TurnReexecutionTime$, so logging lengthens the stride by
   $\sqrt{\TurnReexecutionTime/\TurnReconstructionTime}$. At steady growth, full snapshots write $(\StrideCheckpointCount+1)/2$
   times delta bytes; deltas pay in replay.}{}

\fcard{Group commit}%
  {\text{wait} = \text{gather} + \text{flush}}%
  {A timer gathers $\CommitBatchRecordCount$ intents into one quorum flush. The first arrival waits
   timer plus flush; the last, as the timer fires, the flush.}%
  {Batching divides storage work (flush$/\CommitBatchRecordCount$ per intent), not latency.}{}
\end{rightcol}

% ===================================================================== BACK: PRACTICE
\backside
\begin{answers}
\ans{Finish}{\val{slowcapquarter} + \val{recoveryquarter} = \val{slowtotalquarter}~s; \val{slowcaphalf} +
  \val{recoveryhalf} = \val{slowtotalhalf}~s; \val{stridequarter} still wins, narrowly. $\CheckpointStride^* =
  \sqrt{\val{slowinside}} \approx$ \val{kslow}: a \val{pausefactor}$\times$ pause doubles it (stride card).}
\ans{Group commit}{(a) First \val{qtimer} + \val{qflush} = \val{qfirst}~ms; last \val{qlast}~ms; work
  \val{qflush} $\div$ \val{qbatch} = \val{qwork}~ms per intent. (b) The gather term: first waits
  \val{qtimertwo} + \val{qflush} = \val{qfirsttwo}~ms; work unchanged. (c) A process crash
  only; power loss can erase an intent whose request left (table).}
\ans{Crash positions}{(a) Not dispatched; a found intent looks like (b), so it is held in doubt.
  (b) In doubt: keep the reservation; reconcile. (c) Still in doubt. (d) No: a replica is not a complete prefix (picture).}
\columnbreak
\ans{Resume}{(a) The new epoch: log, gate and resource reject the old one (lease expiry alone
  cannot). (b) No: reconcile by operation identity first (resume idea). (c) Prefill, a floor of about \val{coldprefill}~s for
  \val{coldtokens} tokens; folding calls no model. Prefix-aware placement weighs that against queueing.}
\ans{Divergence}{(a) Harness regression; context assembly.
  (b) That the omission caused the claim. (c) A counterfactual branch restoring the note (table).}
\end{answers}

\begin{leftcol}
\sheetsec{Worked example}
\begin{practicebox}\boxtitle{Which checkpoint stride on this crash trace?}
A \val{turns}-turn repair crashes after turns \val{crashA} and \val{crashB}. A checkpoint pauses
\val{pause}~s (also taken at the last turn); a crash costs \val{restore}~s plus \val{rebuild}~s per lost turn.
Stride \val{strideten}, \val{stridequarter} or \val{stridehalf}?
\par\vspace{5pt}
{\centering
\begin{tikzpicture}[sheet]
  \pgfmathsetmacro\sc{16}
  \foreach \lab/\c/\r/\t/\y/\sty in {\val{strideten} turns/\val{captureten}/\val{recoveryten}/\val{totalten}/32/muted,
      \val{stridequarter} turns/\val{capturequarter}/\val{recoveryquarter}/\val{totalquarter}/16/accent,
      \val{stridehalf} turns/\val{capturehalf}/\val{recoveryhalf}/\val{totalhalf}/0/muted}{
    \node[lab, anchor=east] at (52,\y+5) {\lab};
    \draw[nbox, draw=\sty] (56,\y) rectangle ({56+\c*\sc},\y+10);
    \draw[nbox, draw=\sty, fill=accentsoft] ({56+\c*\sc},\y) rectangle ({56+\t*\sc},\y+10);
    \node[note, anchor=west, text=\sty] at ({60+\t*\sc},\y+5) {\t{} s};}
  \node[note, anchor=south] at ({56+\val{captureten}*\sc/2},43) {capture};
  \node[note, anchor=south] at ({56+(\val{captureten}+\val{totalten})*\sc/2},43) {recovery};
\end{tikzpicture}\par}\vspace{4pt}
\begin{steps}
\item Capture: \val{countten}, \val{countquarter}, \val{counthalf} checkpoints $\times$ \val{pause}~s =
  \val{captureten}, \val{capturequarter}, \val{capturehalf}~s.
\item Recovery: \val{ncrash} $\times$ \val{restore}~s + \val{rebuild}~s $\times$ (\val{lostAten}+\val{lostBten},
  \val{lostAquarter}+\val{lostBquarter}, \val{lostAhalf}+\val{lostBhalf}) = \val{recoveryten}, \val{recoveryquarter}, \val{recoveryhalf}~s.
\item Totals \val{totalten}, \val{totalquarter}, \val{totalhalf}~s: \res{\val{best} turns wins}.
\item {\raggedright Balance: $\TurnFailureRate = \val{ncrash}/\val{turns} = \val{rate}$, so
  $\CheckpointStride^* = \sqrt{2 \times \val{pause} / (\val{rate} \times \val{rebuild})} = \sqrt{\val{balanceinside}}
  \approx$ \res{\val{klog} turns}. Unlogged (\val{reexec}~s per turn): $\sqrt{2 \times \val{pause} / (\val{rate} \times \val{reexec})} =
  \sqrt{\val{execinside}} \approx$ \val{kexec}.\par}
\end{steps}
\takeaway{A logged crash costs reconstruction, not re-execution: the stride grows \val{gain}$\times$.}
\end{practicebox}

\vspace{10pt}
\sheetsec{Your turn: finish it}
\begin{practicebox}\boxtitle{The same trace with a slower capture}
\assumed{} each checkpoint now pauses \val{slowpause}~s.\par\vspace{4pt}
\textbf{Given:} stride \val{strideten}: \val{countten} $\times$ \val{slowpause} + \val{recoveryten} = \val{slowtotalten}~s.\par
\textbf{You:} totals at \val{stridequarter} and \val{stridehalf}, the winner, the new balance stride.\par\vspace{30pt}
\end{practicebox}
\end{leftcol}%
\begin{rightcol}
\sheetsec{Common mistakes}
\mistake{``To reproduce the failure, rerun it at temperature zero.''}%
  {Serving is not numerically deterministic. At \val{rerunq}\% per call, a \val{rerunh}-turn run matches
   $\val{rerunqfrac}^{\val{rerunh}}$ = \val{rerunall}\% of the time. Replay the log.}{}
\mistake{``Checkpoint after every turn to be safe.''}%
  {Logging makes a crash cheap: \val{everyturn}~s against \val{totalquarter}~s
   at \val{best} turns on the worked trace.}{}
\mistake{``Prune the log to save storage.''}%
  {Keeping every recovery input \val{retmonths} months costs \val{retcents} cents, \val{retshare}\% of a
   \$\val{retattempt} attempt. Pruning those bytes forfeits recovery.}{}

\vspace{2pt}
\sheetsec{Check yourself}
\begin{qs}
\item \assumed{} \val{qtimer}~ms timer, \val{qflush}~ms flush, \val{qbatch}-intent batches.
  (a) First and last waits; work per intent? (b) Timer doubled, batch fixed: what grows? (c) Ack from the
  page cache: survives which failures?
\item Recovery states: a crash (a) after the intent append, before its ack: what happened; what does recovery conclude on finding the intent?
  (b) after it, before any result; (c) after a durable timeout. (d) A replica lacks the intent: re-dispatch?
\item A replacement owns the repair; its test is in doubt. The old worker, only partitioned,
  appends again. (a) What stops it? (b) May the new owner re-dispatch the test? (c) On an
  idle worker with no cached prefix, which term must it add?
\item After a harness change, replay shows the next model input dropping
  ``regression suite skipped''; the model then claims success. (a) Which mode; what is localized? (b) What is unproven? (c) How to test the cause?
\end{qs}
\end{rightcol}
\end{document}
